\section{評価方法}

\ref{研究課題と仮説}節で述べた仮説を検証するため，実験を実施するうえで以下の4つの指標について評価を実施する．

\vskip.5\baselineskip

\textbf{自己開示レベル:}自己開示のレベルの推移を明らかにするため，システムログより，各フェーズでの選曲履歴を取得する．
\vskip.5\baselineskip

\textbf{受容度:}仮説1, 2を検証するため，受容性の評価は，被開示者側と開示者側の両面から測定を行う．表\ref{実験中のアンケート項目}に，実験中のアンケート項目を示す．被開示者に対して，相手の自己開示をどの程度理解し，受け入れることができたかについて，表\ref{実験中のアンケート項目}のEQ1, EQ2を使用する．また，開示者側が知覚した相手からの受容についても取得することとして表\ref{実験中のアンケート項目}のEQ5, EQ6を使用する．これにより，段階的な自己開示が受容性に与える効果を，定量的に分析することができる．

また，実験終了後に相手の自己開示に対する印象や，理解が深まった点について表\ref{実験後のアンケート項目}のアンケートを行う．具体的には，相手に対する理解度，相手の音楽の好みに対する理解度，今後も相手と音楽を通じた会話をしたいと思うかを評価項目として設定する．
\vskip.5\baselineskip

\textbf{曲への興味:}仮説2を検証するため，被開示者に対し，選択された楽曲を知っていたかどうか（EQ3）と，相手の曲に興味を持てたかどうか（EQ4）という指標で相手の選曲に対しての評価を行う．これにより，音楽への興味度と受容性の関係性を分析する．
\vskip.5\baselineskip

\textbf{被験者の意見:}実験後のヒアリング項目を表\ref{実験後のヒアリング項目}に示す．選曲の意識，セッション全体の雰囲気，フェーズ前半と後半での違いなどについて，被験者からの主観的な評価を収集する．

\vskip\baselineskip

アンケートは，一部はい，いいえの2値や複数選択の回答があるが，明記していないものはすべて5段階評価（1:全くそう思わない ～ 5:とてもそう思う）を使用している．

\begin{table}[H]
  \caption{実験中のアンケート項目（EQ）}
  \label{実験中のアンケート項目}
  \centering
  \begin{tabular}{|l|p{0.85\textwidth}|}
    \hline
     & \textbf{質問内容} \\
    \hhline{|=|=|}
    EQ1 & 流した曲について楽しく話ができた \\
    \hline
    EQ2 & 相手に自分の思いが理解してもらえたと感じた \\
    \hline
    EQ3 & 相手の曲を知っていた（はい，いいえ） \\
    \hline
    EQ4 & 相手の曲に興味を持てた \\
    \hline
    EQ5 & 聴いた曲について相手と楽しく話ができた \\
    \hline
    EQ6 & 相手の曲に関する思いや経験が理解できた \\
    \hline
  \end{tabular}
\end{table}

\begin{table}[H]
  \caption{実験後のアンケート項目（PQ）}
  \label{実験後のアンケート項目}
  \centering
  \begin{tabular}{|l|p{0.85\textwidth}|}
    \hline
    \multicolumn{2}{|c|}{\textbf{音楽に関する基本情報}} \\
    \hhline{|=|=|}
    PQ1 & 普段聴いている音楽ジャンルの傾向（複数選択） \\
    \hline
    PQ2 & 相手が知らない曲を話すことに抵抗がある \\
    \hline
    PQ3 & 知らない曲でも楽しむことができる \\
    \hhline{|=|=|}
    \multicolumn{2}{|c|}{\textbf{自己開示と受容について}} \\
    \hhline{|=|=|}
    PQ4 & 自分の内面についてよく話すことができた \\
    \hline
    PQ5 & 相手のことをより深く理解することができた \\
    \hline
    PQ6 & 実験前より相手の音楽の好みについて理解できた \\
    \hline
    PQ7 & 今後も相手と音楽を通じた会話をしてみたい \\
    \hline
  \end{tabular}
\end{table}

\begin{table}[H]
  \caption{実験後のヒアリング項目（HQ）}
  \label{実験後のヒアリング項目}
  \centering
  \begin{tabular}{|l|p{0.85\textwidth}|}
    \hline
     & \textbf{質問内容} \\
    \hhline{|=|=|}
    HQ1 & セッション全体を通して，楽しむことができましたか．また，何が楽しかったですか \\
    \hline
    HQ2 & どのような基準で楽曲を選択しましたか \\
    \hline
    HQ3 & フェーズが進むにつれて，楽曲選択の意識に変化はありましたか \\
    \hline
    HQ4 & 音楽を共有しながら会話をすることについて，どのように感じましたか \\
    \hline
    HQ5 & 特に会話のしやすさを感じたフェーズ，または会話のしづらさを感じたフェーズはありましたか．その理由もお聞かせください \\
    \hline
    HQ6 & 前半（1-4回目）と後半（5-8回目）で，相手との会話に変化を感じましたか \\
    \hline
    HQ7 & システムからの推薦された曲について，話しやすかったですか \\
       &（推薦あり群のみ）\\
    \hline
  \end{tabular}
\end{table}
